\documentclass{article}


\title{Ansätze und Überlegungen}
\author{Bjarne Jehn}
\date{2026}

\begin{document}
	
\maketitle

\begin{itemize}
	\item greedy idee: für ein level lösen, dann für ein weiteres, dann für ein weiteres - wird nicht optimal werden aber eventuell besser als nur für ein level
	\item suchräume verkleinern: sind die parameter tatsächlich non-decreasing? Wenn ja, Begründen und eventuell ausnutzen
	\item locate case modellieren - einfach nur $f_{avg}/2$ durch $f_{avg}$ ersetzten?
	\item $p_M$ bzw. $p_K$ als 6-dimensionale funktion mit parametern $(m^1,m^2,m^3,t^1,t^2,t^3)$ betrachten. oder vllt erst mal als 3-dimensionale mit fixen thresholds um ein gefühl dafür zu bekommen
\end{itemize}
	
	
\end{document}